\section{核验锚点簇：洪武：城池接收、官署与北方压力（1368--1398）}
\label{register:1368-1398_hongwu}
开国编年不能把诏令、攻城与地方接收压成同一种胜利。本簇让每条记录回到其所说的地点或空间尺度，并保留军报、官书和地方社会之间无法由一份材料填平的距离。

以下为读者可跳转的首次描写。它们是按王朝年序编排的证据性短条，而非把账本字段伪装成同一种长篇叙事：每项只在所列材料可承受的范围内给出行动者、空间、后果与限度。

\begin{description}
  \item[\eventanchor{EM-1368-DADU}{1368（洪武元年）八月：徐达军进入大都并改称北平}] 明与北元；朝廷与官员、军队与卫所；跨区域行政、资源或军政网络；不假定同步执行。明军自山东、河南北上且元顺帝已北撤；大都守备难获外援 元廷失去旧都和华北行政中心；明须经营北平并面对草原残余政权 材料：《太祖实录》洪武元年八月条；《明史·太祖本纪》；《高丽史·恭愍王世家》二十一年条。限度：日期可靠；明廷军报不独证城内损失和居民去留
  \item[\eventanchor{EM-1368-GUOZIJIAN}{1368（洪武元年）九月：在南京设国子学以培训中央官员子弟}] 明；朝廷与官员、宗教与文化行动者；跨区域行政、资源或军政网络；不假定同步执行。新朝须为六部和地方官体系补充识字行政人才并重建礼制教育 南京国子监成为明初官学枢纽；科举与荐举的关系仍待后续制度调整 材料：《太祖实录》洪武元年九月条；《明会典·国子监》；《明史·选举志》。限度：日期较可靠；学校建置不等于生员规模或教学成效
  \item[\eventanchor{ML-1368-001}{1368（洪武元年）一月18日：红巾起义：胡美攻占福州}] 明与地方反抗、流动作战集团；朝廷与官员、军队与卫所、地方社会与精英、流动与反抗集团；地方或省域；未具城镇者不补造路线。前期边防、地方武装或跨区域资源调动的压力推动了该行动。 它改变了后续调兵、补给或地方秩序的条件；战果和伤亡须分开核对。 材料：《太祖实录》洪武元年条；《明史·兵志》及相关列传；边镇、地方志或对方材料。限度：译写候选以编年材料定位；实录记载反映朝廷选择，崇祯期须另以奏疏、地方及敌对政权材料交叉。
  \item[\eventanchor{ML-1368-002}{1368（洪武元年）一月23日：朱元璋自立为明朝洪武皇帝（注意明清使用年号而不是庙名）}] 明、后金（清）与辽东诸势力；朝廷与官员、边地政权与社群；边地接触区；不把族名或部名当固定疆界。继承、官僚协调或皇权运作的压力使问题进入朝廷程序。 它影响后续人事与政策议程；动机和派系标签不能由单一叙事裁定。 材料：《太祖实录》洪武元年条；《明史·本纪》及相关列传；诏令、奏疏或会典版本。限度：译写候选以编年材料定位；实录记载反映朝廷选择，崇祯期须另以奏疏、地方及敌对政权材料交叉。
  \item[\eventanchor{ML-1368-003}{1368（洪武元年）二月17日：明军攻克福建，俘获陈有定，并处决}] 明；朝廷与官员、军队与卫所、地方社会与精英；地方或省域；未具城镇者不补造路线。前期边防、地方武装或跨区域资源调动的压力推动了该行动。 它改变了后续调兵、补给或地方秩序的条件；战果和伤亡须分开核对。 材料：《太祖实录》洪武元年条；《明史·兵志》及相关列传；边镇、地方志或对方材料。限度：译写候选以编年材料定位；实录记载反映朝廷选择，崇祯期须另以奏疏、地方及敌对政权材料交叉。
  \item[\eventanchor{ML-1368-004}{1368（洪武元年）三月1日：明军攻克山东}] 明；朝廷与官员、军队与卫所、地方社会与精英；地方或省域；未具城镇者不补造路线。前期边防、地方武装或跨区域资源调动的压力推动了该行动。 它改变了后续调兵、补给或地方秩序的条件；战果和伤亡须分开核对。 材料：《太祖实录》洪武元年条；《明史·兵志》及相关列传；边镇、地方志或对方材料。限度：译写候选以编年材料定位；实录记载反映朝廷选择，崇祯期须另以奏疏、地方及敌对政权材料交叉。
  \item[\eventanchor{ML-1368-005}{1368（洪武元年）四月16日：明军攻克开封}] 明；朝廷与官员、军队与卫所；跨区域行政、资源或军政网络；不假定同步执行。前期边防、地方武装或跨区域资源调动的压力推动了该行动。 它改变了后续调兵、补给或地方秩序的条件；战果和伤亡须分开核对。 材料：《太祖实录》洪武元年条；《明史·兵志》及相关列传；边镇、地方志或对方材料。限度：译写候选以编年材料定位；实录记载反映朝廷选择，崇祯期须另以奏疏、地方及敌对政权材料交叉。
  \item[\eventanchor{ML-1368-006}{1368（洪武元年）四月18日：明军到达广州并接受何震的投降}] 明；朝廷与官员、军队与卫所；跨区域行政、资源或军政网络；不假定同步执行。前期边防、地方武装或跨区域资源调动的压力推动了该行动。 它改变了后续调兵、补给或地方秩序的条件；战果和伤亡须分开核对。 材料：《太祖实录》洪武元年条；《明史·兵志》及相关列传；边镇、地方志或对方材料。限度：译写候选以编年材料定位；实录记载反映朝廷选择，崇祯期须另以奏疏、地方及敌对政权材料交叉。
  \item[\eventanchor{ML-1368-007}{1368（洪武元年）四月25日：明军击败阔克帖木儿并占领洛阳}] 明；朝廷与官员、军队与卫所；跨区域行政、资源或军政网络；不假定同步执行。前期边防、地方武装或跨区域资源调动的压力推动了该行动。 它改变了后续调兵、补给或地方秩序的条件；战果和伤亡须分开核对。 材料：《太祖实录》洪武元年条；《明史·兵志》及相关列传；边镇、地方志或对方材料。限度：译写候选以编年材料定位；实录记载反映朝廷选择，崇祯期须另以奏疏、地方及敌对政权材料交叉。
  \item[\eventanchor{ML-1368-009}{1368（洪武元年）七月：明军攻克广西}] 明；朝廷与官员、军队与卫所；跨区域行政、资源或军政网络；不假定同步执行。前期边防、地方武装或跨区域资源调动的压力推动了该行动。 它改变了后续调兵、补给或地方秩序的条件；战果和伤亡须分开核对。 材料：《太祖实录》洪武元年条；《明史·兵志》及相关列传；边镇、地方志或对方材料。限度：译写候选以编年材料定位；实录记载反映朝廷选择，崇祯期须另以奏疏、地方及敌对政权材料交叉。
  \item[\eventanchor{ML-1368-010}{1368（洪武元年）九月20日：明军占领大都（改称北平），元朝逃往内蒙古；元朝就这样结束了}] 明与蒙古诸部；朝廷与官员、军队与卫所、边地政权与社群；边地接触区；不把族名或部名当固定疆界。前期边防、地方武装或跨区域资源调动的压力推动了该行动。 它改变了后续调兵、补给或地方秩序的条件；战果和伤亡须分开核对。 材料：《太祖实录》洪武元年条；《明史·兵志》及相关列传；边镇、地方志或对方材料。限度：译写候选以编年材料定位；实录记载反映朝廷选择，崇祯期须另以奏疏、地方及敌对政权材料交叉。
  \item[\eventanchor{ML-1368-011}{1368（洪武元年）十一月：明军攻克保定}] 明；朝廷与官员、军队与卫所；跨区域行政、资源或军政网络；不假定同步执行。前期边防、地方武装或跨区域资源调动的压力推动了该行动。 它改变了后续调兵、补给或地方秩序的条件；战果和伤亡须分开核对。 材料：《太祖实录》洪武元年条；《明史·兵志》及相关列传；边镇、地方志或对方材料。限度：译写候选以编年材料定位；实录记载反映朝廷选择，崇祯期须另以奏疏、地方及敌对政权材料交叉。
  \item[\eventanchor{ML-1368-012}{1368（洪武元年）十二月26日：明军攻克赵州}] 明；朝廷与官员、军队与卫所；跨区域行政、资源或军政网络；不假定同步执行。前期边防、地方武装或跨区域资源调动的压力推动了该行动。 它改变了后续调兵、补给或地方秩序的条件；战果和伤亡须分开核对。 材料：《太祖实录》洪武元年条；《明史·兵志》及相关列传；边镇、地方志或对方材料。限度：译写候选以编年材料定位；实录记载反映朝廷选择，崇祯期须另以奏疏、地方及敌对政权材料交叉。
  \item[\eventanchor{ML-1368-013}{1368（洪武元年）十二月：明军攻克平定}] 明；朝廷与官员、军队与卫所；跨区域行政、资源或军政网络；不假定同步执行。前期边防、地方武装或跨区域资源调动的压力推动了该行动。 它改变了后续调兵、补给或地方秩序的条件；战果和伤亡须分开核对。 材料：《太祖实录》洪武元年条；《明史·兵志》及相关列传；边镇、地方志或对方材料。限度：译写候选以编年材料定位；实录记载反映朝廷选择，崇祯期须另以奏疏、地方及敌对政权材料交叉。
  \item[\eventanchor{ML-1368-014}{1368（洪武元年）：卧虎炮为明军所用。}] 明；朝廷与官员、军队与卫所；跨区域行政、资源或军政网络；不假定同步执行。制度、教育或知识传播的需要推动了相关行动或文本形成。 它提供制度和传播史的锚点，不能据此推断社会整体接受。 材料：《太祖实录》洪武元年条；《明史·选举志》或《艺文志》；文集、碑刻或书目材料。限度：译写候选以编年材料定位；实录记载反映朝廷选择，崇祯期须另以奏疏、地方及敌对政权材料交叉。
  \item[\eventanchor{ML-1368-015}{1368（洪武元年）：国子监创建}] 明；朝廷与官员；跨区域行政、资源或军政网络；不假定同步执行。继承、官僚协调或皇权运作的压力使问题进入朝廷程序。 它影响后续人事与政策议程；动机和派系标签不能由单一叙事裁定。 材料：《太祖实录》洪武元年条；《明史·本纪》及相关列传；诏令、奏疏或会典版本。限度：译写候选以编年材料定位；实录记载反映朝廷选择，崇祯期须另以奏疏、地方及敌对政权材料交叉。
  \item[\eventanchor{EM-1369-SHANGDU}{1369（洪武二年）五月至八月：明军北进上都、应昌一线；元顺帝在北遁途中去世}] 明与北元；朝廷与官员、军队与卫所；跨区域行政、资源或军政网络；不假定同步执行。大都失守后北元仍试图依托上都、应昌和草原骑兵维持王朝中心 明军未能消灭草原政治力量；北边战争由夺城转为长期机动与互市问题 材料：《太祖实录》洪武二年五月、八月条；《明史·北虏传》；《高丽史》北元消息条。限度：日期与战事序列较可靠；元顺帝死地和撤退细节有异说
  \item[\eventanchor{ML-1369-001}{1369（洪武二年）一月9日：明军攻克太原}] 明；朝廷与官员、军队与卫所；跨区域行政、资源或军政网络；不假定同步执行。前期边防、地方武装或跨区域资源调动的压力推动了该行动。 它改变了后续调兵、补给或地方秩序的条件；战果和伤亡须分开核对。 材料：《太祖实录》洪武二年条；《明史·兵志》及相关列传；边镇、地方志或对方材料。限度：译写候选以编年材料定位；实录记载反映朝廷选择，崇祯期须另以奏疏、地方及敌对政权材料交叉。
  \item[\eventanchor{ML-1369-002}{1369（洪武二年）三月3日：明军攻克大同}] 明；朝廷与官员、军队与卫所；跨区域行政、资源或军政网络；不假定同步执行。前期边防、地方武装或跨区域资源调动的压力推动了该行动。 它改变了后续调兵、补给或地方秩序的条件；战果和伤亡须分开核对。 材料：《太祖实录》洪武二年条；《明史·兵志》及相关列传；边镇、地方志或对方材料。限度：译写候选以编年材料定位；实录记载反映朝廷选择，崇祯期须另以奏疏、地方及敌对政权材料交叉。
  \item[\eventanchor{ML-1369-003}{1369（洪武二年）三月：宋濂、王邑开始编修《元史》}] 明；朝廷与官员；跨区域行政、资源或军政网络；不假定同步执行。继承、官僚协调或皇权运作的压力使问题进入朝廷程序。 它影响后续人事与政策议程；动机和派系标签不能由单一叙事裁定。 材料：《太祖实录》洪武二年条；《明史·本纪》及相关列传；诏令、奏疏或会典版本。限度：译写候选以编年材料定位；实录记载反映朝廷选择，崇祯期须另以奏疏、地方及敌对政权材料交叉。
  \item[\eventanchor{ML-1369-004}{1369（洪武二年）四月18日：明军攻克山西，李思齐逃往临洮}] 明；朝廷与官员、军队与卫所；跨区域行政、资源或军政网络；不假定同步执行。前期边防、地方武装或跨区域资源调动的压力推动了该行动。 它改变了后续调兵、补给或地方秩序的条件；战果和伤亡须分开核对。 材料：《太祖实录》洪武二年条；《明史·兵志》及相关列传；边镇、地方志或对方材料。限度：译写候选以编年材料定位；实录记载反映朝廷选择，崇祯期须另以奏疏、地方及敌对政权材料交叉。
  \item[\eventanchor{ML-1369-005}{1369（洪武二年）五月21日：李思齐向明军投降}] 明；朝廷与官员、军队与卫所；跨区域行政、资源或军政网络；不假定同步执行。前期边防、地方武装或跨区域资源调动的压力推动了该行动。 它改变了后续调兵、补给或地方秩序的条件；战果和伤亡须分开核对。 材料：《太祖实录》洪武二年条；《明史·兵志》及相关列传；边镇、地方志或对方材料。限度：译写候选以编年材料定位；实录记载反映朝廷选择，崇祯期须另以奏疏、地方及敌对政权材料交叉。
  \item[\eventanchor{ML-1369-006}{1369（洪武二年）五月23日：明军攻克兰州}] 明；朝廷与官员、军队与卫所；跨区域行政、资源或军政网络；不假定同步执行。前期边防、地方武装或跨区域资源调动的压力推动了该行动。 它改变了后续调兵、补给或地方秩序的条件；战果和伤亡须分开核对。 材料：《太祖实录》洪武二年条；《明史·兵志》及相关列传；边镇、地方志或对方材料。限度：译写候选以编年材料定位；实录记载反映朝廷选择，崇祯期须另以奏疏、地方及敌对政权材料交叉。
  \item[\eventanchor{ML-1369-007}{1369（洪武二年）六月8日：明军攻克平凉}] 明；朝廷与官员、军队与卫所；跨区域行政、资源或军政网络；不假定同步执行。前期边防、地方武装或跨区域资源调动的压力推动了该行动。 它改变了后续调兵、补给或地方秩序的条件；战果和伤亡须分开核对。 材料：《太祖实录》洪武二年条；《明史·兵志》及相关列传；边镇、地方志或对方材料。限度：译写候选以编年材料定位；实录记载反映朝廷选择，崇祯期须另以奏疏、地方及敌对政权材料交叉。
  \item[\eventanchor{ML-1369-008}{1369（洪武二年）七月20日：明军攻克上都}] 明；朝廷与官员、军队与卫所；跨区域行政、资源或军政网络；不假定同步执行。前期边防、地方武装或跨区域资源调动的压力推动了该行动。 它改变了后续调兵、补给或地方秩序的条件；战果和伤亡须分开核对。 材料：《太祖实录》洪武二年条；《明史·兵志》及相关列传；边镇、地方志或对方材料。限度：译写候选以编年材料定位；实录记载反映朝廷选择，崇祯期须另以奏疏、地方及敌对政权材料交叉。
  \item[\eventanchor{ML-1369-009}{1369（洪武二年）九月22日：明军攻克庆阳}] 明、后金（清）与辽东诸势力；朝廷与官员、军队与卫所、边地政权与社群；边地接触区；不把族名或部名当固定疆界。前期边防、地方武装或跨区域资源调动的压力推动了该行动。 它改变了后续调兵、补给或地方秩序的条件；战果和伤亡须分开核对。 材料：《太祖实录》洪武二年条；《明史·兵志》及相关列传；边镇、地方志或对方材料。限度：译写候选以编年材料定位；实录记载反映朝廷选择，崇祯期须另以奏疏、地方及敌对政权材料交叉。
  \item[\eventanchor{EM-1370-PRINCELY-ENFEOFFMENT}{1370（洪武三年）四月：太祖大封诸子为王并配置藩国}] 明；朝廷与官员；跨区域行政、资源或军政网络；不假定同步执行。建国后皇室需要可继承的资源分配与边地防卫网络 诸王成为地方军政力量；中央与藩王的权限张力为建文时期冲突埋下制度条件 材料：《太祖实录》洪武三年四月条；《明史·诸王传》；《明会典·宗藩》。限度：日期可靠；封册可证名义授封而不能直接证明各王实际兵权
  \item[\eventanchor{EM-1370-SHENER-BATTLE}{1370（洪武三年）六月：徐达在沈儿峪击败扩廓帖木儿部}] 明与北元；朝廷与官员；跨区域行政、资源或军政网络；不假定同步执行。明廷欲趁北元重组未定追击其主力；徐达军依托山西、陕西粮道北进 扩廓部受挫但未被歼灭；明廷随后仍需以多路远征争夺草原主动权 材料：《太祖实录》洪武三年六月条；《明史·徐达传》与《北虏传》；17世纪《蒙古源流》北元汗统叙事。限度：战役发生可靠；双方兵数和斩获主要来自明方战报
  \item[\eventanchor{ML-1370-001}{1370（洪武三年）一月：阔克帖木儿围攻兰州但未能攻克}] 明；朝廷与官员、军队与卫所；跨区域行政、资源或军政网络；不假定同步执行。前期边防、地方武装或跨区域资源调动的压力推动了该行动。 它改变了后续调兵、补给或地方秩序的条件；战果和伤亡须分开核对。 材料：《太祖实录》洪武三年条；《明史·兵志》及相关列传；边镇、地方志或对方材料。限度：译写候选以编年材料定位；实录记载反映朝廷选择，崇祯期须另以奏疏、地方及敌对政权材料交叉。
  \item[\eventanchor{ML-1370-002}{1370（洪武三年）：明军在巩昌击败了可克帖木儿，但未能俘获他}] 明；朝廷与官员、军队与卫所；跨区域行政、资源或军政网络；不假定同步执行。前期边防、地方武装或跨区域资源调动的压力推动了该行动。 它改变了后续调兵、补给或地方秩序的条件；战果和伤亡须分开核对。 材料：《太祖实录》洪武三年条；《明史·兵志》及相关列传；边镇、地方志或对方材料。限度：译写候选以编年材料定位；实录记载反映朝廷选择，崇祯期须另以奏疏、地方及敌对政权材料交叉。
  \item[\eventanchor{ML-1370-003}{1370（洪武三年）六月5日：洪武皇帝授权明第一次科举考试}] 明；朝廷与官员、宗教与文化行动者；跨区域行政、资源或军政网络；不假定同步执行。制度、教育或知识传播的需要推动了相关行动或文本形成。 它提供制度和传播史的锚点，不能据此推断社会整体接受。 材料：《太祖实录》洪武三年条；《明史·选举志》或《艺文志》；文集、碑刻或书目材料。限度：译写候选以编年材料定位；实录记载反映朝廷选择，崇祯期须另以奏疏、地方及敌对政权材料交叉。
  \item[\eventanchor{ML-1370-004}{1370（洪武三年）六月10日：明军攻克颍昌}] 明；朝廷与官员、军队与卫所；跨区域行政、资源或军政网络；不假定同步执行。前期边防、地方武装或跨区域资源调动的压力推动了该行动。 它改变了后续调兵、补给或地方秩序的条件；战果和伤亡须分开核对。 材料：《太祖实录》洪武三年条；《明史·兵志》及相关列传；边镇、地方志或对方材料。限度：译写候选以编年材料定位；实录记载反映朝廷选择，崇祯期须另以奏疏、地方及敌对政权材料交叉。
  \item[\eventanchor{ML-1370-005}{1370（洪武三年）六月：洪武帝取缔白莲教和摩尼教}] 明；朝廷与官员、流动与反抗集团、宗教与文化行动者；跨区域行政、资源或军政网络；不假定同步执行。继承、官僚协调或皇权运作的压力使问题进入朝廷程序。 它影响后续人事与政策议程；动机和派系标签不能由单一叙事裁定。 材料：《太祖实录》洪武三年条；《明史·本纪》及相关列传；诏令、奏疏或会典版本。限度：译写候选以编年材料定位；实录记载反映朝廷选择，崇祯期须另以奏疏、地方及敌对政权材料交叉。
  \item[\eventanchor{ML-1370-006}{1370（洪武三年）：明代使用火药来增强地雷的爆炸力。}] 明；朝廷与官员；跨区域行政、资源或军政网络；不假定同步执行。财政征解、交通工程或货币支付的压力使相关措施进入政务。 其法定安排不等于地方执行，后续须以地域材料追踪负担与成效。 材料：《太祖实录》洪武三年条；《明会典》相关条目；地方志、黄册/契约/河工材料。限度：译写候选以编年材料定位；实录记载反映朝廷选择，崇祯期须另以奏疏、地方及敌对政权材料交叉。
  \item[\eventanchor{ML-1370-007}{1370（洪武三年）：明代大炮的弹丸由石弹过渡为铁弹。}] 明；朝廷与官员；跨区域行政、资源或军政网络；不假定同步执行。制度、教育或知识传播的需要推动了相关行动或文本形成。 它提供制度和传播史的锚点，不能据此推断社会整体接受。 材料：《太祖实录》洪武三年条；《明史·选举志》或《艺文志》；文集、碑刻或书目材料。限度：译写候选以编年材料定位；实录记载反映朝廷选择，崇祯期须另以奏疏、地方及敌对政权材料交叉。
  \item[\eventanchor{EM-1371-HAIJIN}{1371（洪武四年）十二月：朝廷申禁滨海居民私自出海并限制民间海外贸易}] 明与东南沿海社会；朝廷与官员、财政与商业行动者；账本可辨空间；地点细节仍须回到所列材料。倭患警报、元末海上武装遗绪与新朝对沿海人口流动的疑虑交织 朝贡贸易被置于官府控制之下；私航、走私和沿海防务成为持续的地方—中央博弈 材料：《太祖实录》洪武四年十二月条；《明史·食货志》；《筹海图编》所录早明海禁材料。限度：诏令日期可靠；禁令不等于沿海航行和交易立即停止
  \item[\eventanchor{EM-1371-MINGXIA}{1371（洪武四年）：傅友德等沿长江入蜀，夏主明昇降明}] 明与明夏；朝廷与官员；跨区域行政、资源或军政网络；不假定同步执行。明夏控制四川但与中原补给和政治盟友隔绝；明军掌握长江中下游军运条件 四川纳入明朝行政与赋役框架；山地边区和土司关系仍须逐步处理 材料：《太祖实录》洪武四年条；《明史·明玉珍传》；《四川通志》明夏遗迹条。限度：年代可靠；行军路线和地方响应不能只据凯旋叙事复原
  \item[\eventanchor{ML-1371-001}{1371（洪武四年）五月18日：明军攻占温州}] 明；朝廷与官员、军队与卫所；跨区域行政、资源或军政网络；不假定同步执行。前期边防、地方武装或跨区域资源调动的压力推动了该行动。 它改变了后续调兵、补给或地方秩序的条件；战果和伤亡须分开核对。 材料：《太祖实录》洪武四年条；《明史·兵志》及相关列传；边镇、地方志或对方材料。限度：译写候选以编年材料定位；实录记载反映朝廷选择，崇祯期须另以奏疏、地方及敌对政权材料交叉。
  \item[\eventanchor{ML-1371-002}{1371（洪武四年）七月：明军攻克杭州}] 明；朝廷与官员、军队与卫所；跨区域行政、资源或军政网络；不假定同步执行。前期边防、地方武装或跨区域资源调动的压力推动了该行动。 它改变了后续调兵、补给或地方秩序的条件；战果和伤亡须分开核对。 材料：《太祖实录》洪武四年条；《明史·兵志》及相关列传；边镇、地方志或对方材料。限度：译写候选以编年材料定位；实录记载反映朝廷选择，崇祯期须另以奏疏、地方及敌对政权材料交叉。
  \item[\eventanchor{ML-1371-003}{1371（洪武四年）八月3日：明盛将四川投降明朝}] 明；朝廷与官员、地方社会与精英；地方或省域；未具城镇者不补造路线。前期边防、地方武装或跨区域资源调动的压力推动了该行动。 它改变了后续调兵、补给或地方秩序的条件；战果和伤亡须分开核对。 材料：《太祖实录》洪武四年条；《明史·兵志》及相关列传；边镇、地方志或对方材料。限度：译写候选以编年材料定位；实录记载反映朝廷选择，崇祯期须另以奏疏、地方及敌对政权材料交叉。
  \item[\eventanchor{ML-1371-004}{1371（洪武四年）：国子监注册学生达3,728人}] 明；朝廷与官员、宗教与文化行动者；跨区域行政、资源或军政网络；不假定同步执行。继承、官僚协调或皇权运作的压力使问题进入朝廷程序。 它影响后续人事与政策议程；动机和派系标签不能由单一叙事裁定。 材料：《太祖实录》洪武四年条；《明史·本纪》及相关列传；诏令、奏疏或会典版本。限度：译写候选以编年材料定位；实录记载反映朝廷选择，崇祯期须另以奏疏、地方及敌对政权材料交叉。
  \item[\eventanchor{EM-1372-LINGBEI-EXPEDITION}{1372（洪武五年）：三路北伐受挫，李文忠、徐达等部未能取得预期战果}] 明与北元；朝廷与官员、军队与卫所；跨区域行政、资源或军政网络；不假定同步执行。明廷试图以多路大军压迫北元；行军距离、草原水草和部队协同限制远征 太祖转向更审慎的北边经营；卫所、边墙和藩王防线的重要性上升 材料：《太祖实录》洪武五年条；《明史·李文忠传》与《北虏传》；17世纪《蒙古源流》北元汗统叙事。限度：战役大势可靠；败因、损失和各路会师情况在记载中不一
  \item[\eventanchor{EM-1372-RYUKYU}{1372（洪武五年）：明使杨载出使琉球并开启朝贡往来}] 明与琉球诸政权；朝廷与官员；账本可辨空间；地点细节仍须回到所列材料。海禁后明廷仍需以官方使节组织可控的海上外交；琉球各势力寻求外部承认 形成持续的册封—贸易通道；朝贡名义与实际航运、商人网络不能混同 材料：《太祖实录》洪武五年条；《明史·琉球传》；《历代宝案》早期入贡文书。限度：明使出航与琉球回应有双边记录；三山内部政治不应视作单一国家
  \item[\eventanchor{ML-1372-002}{1372（洪武五年）：明军在喀喇昆仑被击溃}] 明；朝廷与官员、军队与卫所；跨区域行政、资源或军政网络；不假定同步执行。前期边防、地方武装或跨区域资源调动的压力推动了该行动。 它改变了后续调兵、补给或地方秩序的条件；战果和伤亡须分开核对。 材料：《太祖实录》洪武五年条；《明史·兵志》及相关列传；边镇、地方志或对方材料。限度：译写候选以编年材料定位；实录记载反映朝廷选择，崇祯期须另以奏疏、地方及敌对政权材料交叉。
  \item[\eventanchor{ML-1372-003}{1372（洪武五年）：明军攻克永昌，攻克居延}] 明；朝廷与官员、军队与卫所；跨区域行政、资源或军政网络；不假定同步执行。前期边防、地方武装或跨区域资源调动的压力推动了该行动。 它改变了后续调兵、补给或地方秩序的条件；战果和伤亡须分开核对。 材料：《太祖实录》洪武五年条；《明史·兵志》及相关列传；边镇、地方志或对方材料。限度：译写候选以编年材料定位；实录记载反映朝廷选择，崇祯期须另以奏疏、地方及敌对政权材料交叉。
  \item[\eventanchor{ML-1372-004}{1372（洪武五年）：国子监注册学生突破1万人}] 明；朝廷与官员、宗教与文化行动者；跨区域行政、资源或军政网络；不假定同步执行。继承、官僚协调或皇权运作的压力使问题进入朝廷程序。 它影响后续人事与政策议程；动机和派系标签不能由单一叙事裁定。 材料：《太祖实录》洪武五年条；《明史·本纪》及相关列传；诏令、奏疏或会典版本。限度：译写候选以编年材料定位；实录记载反映朝廷选择，崇祯期须另以奏疏、地方及敌对政权材料交叉。
  \item[\eventanchor{ML-1372-005}{1372（洪武五年）：明代出现了专门为海军使用的大炮。}] 明；朝廷与官员、军队与卫所；账本可辨空间；地点细节仍须回到所列材料。制度、教育或知识传播的需要推动了相关行动或文本形成。 它提供制度和传播史的锚点，不能据此推断社会整体接受。 材料：《太祖实录》洪武五年条；《明史·选举志》或《艺文志》；文集、碑刻或书目材料。限度：译写候选以编年材料定位；实录记载反映朝廷选择，崇祯期须另以奏疏、地方及敌对政权材料交叉。
  \item[\eventanchor{EM-1373-EXAMINATION-SUSPENDED}{1373（洪武六年）：太祖停罢科举，改重荐举与学校考核}] 明；朝廷与官员、宗教与文化行动者；跨区域行政、资源或军政网络；不假定同步执行。首次科举所取士人与新朝行政需求之间发生张力；皇帝偏好直接控制人才来源 荐举和学校渠道暂时扩张；科举何时、如何恢复成为中央选官制度的后续问题 材料：《太祖实录》洪武六年条；《明史·选举志》；《明会典·礼部》科举沿革。限度：政策日期较可靠；官方对取士质量的评语带有皇帝立场
  \item[\eventanchor{ML-1373-001}{1373（洪武六年）三月：洪武皇帝暂停科举考试}] 明；朝廷与官员、宗教与文化行动者；跨区域行政、资源或军政网络；不假定同步执行。制度、教育或知识传播的需要推动了相关行动或文本形成。 它提供制度和传播史的锚点，不能据此推断社会整体接受。 材料：《太祖实录》洪武六年条；《明史·选举志》或《艺文志》；文集、碑刻或书目材料。限度：译写候选以编年材料定位；实录记载反映朝廷选择，崇祯期须另以奏疏、地方及敌对政权材料交叉。
  \item[\eventanchor{ML-1373-002}{1373（洪武六年）十一月29日：明军在怀柔击败可克帖木儿}] 明；朝廷与官员、军队与卫所；跨区域行政、资源或军政网络；不假定同步执行。前期边防、地方武装或跨区域资源调动的压力推动了该行动。 它改变了后续调兵、补给或地方秩序的条件；战果和伤亡须分开核对。 材料：《太祖实录》洪武六年条；《明史·兵志》及相关列传；边镇、地方志或对方材料。限度：译写候选以编年材料定位；实录记载反映朝廷选择，崇祯期须另以奏疏、地方及敌对政权材料交叉。
  \item[\eventanchor{ML-1373-003}{1373（洪武六年）：洪武皇帝将高丽的进贡限制为每三年一次}] 明；朝廷与官员；跨区域行政、资源或军政网络；不假定同步执行。朝贡名义、边境安全与港口贸易的利益使交涉持续发生。 它为后续册封、互市或海防安排留下文书与跨政权比较线索。 材料：《太祖实录》洪武六年条；《明史·外国传》；《朝鲜王朝实录》、琉球《历代宝案》或海外档案。限度：译写候选以编年材料定位；实录记载反映朝廷选择，崇祯期须另以奏疏、地方及敌对政权材料交叉。
  \item[\eventanchor{ML-1373-004}{1373（洪武六年）：明朝官员制定了中国历史上第一部《宅法》}] 明；朝廷与官员；跨区域行政、资源或军政网络；不假定同步执行。继承、官僚协调或皇权运作的压力使问题进入朝廷程序。 它影响后续人事与政策议程；动机和派系标签不能由单一叙事裁定。 材料：《太祖实录》洪武六年条；《明史·本纪》及相关列传；诏令、奏疏或会典版本。限度：译写候选以编年材料定位；实录记载反映朝廷选择，崇祯期须另以奏疏、地方及敌对政权材料交叉。
  \item[\eventanchor{ML-1374-001}{1374（洪武七年）：朝廷围绕卫所军户的编籍、屯田与调发整理本年军政文书，作为明初军事接收的可核查节点。}] 明；朝廷与官员、军队与卫所、财政与商业行动者；跨区域行政、资源或军政网络；不假定同步执行。前期边防、地方武装或跨区域资源调动的压力推动了该行动。 它改变了后续调兵、补给或地方秩序的条件；战果和伤亡须分开核对。 材料：《太祖实录》洪武七年条；《明史·兵志》及相关列传；边镇、地方志或对方材料。限度：桥接条目只标示可回查的年度问题与材料链；实录有中央选择性，崇祯期尤其需要奏疏、地方、朝鲜及满文材料交叉。
  \item[\eventanchor{ML-1374-002}{1374（洪武七年）：北元诸部、东北和西南的边报、招抚与防御命令持续进入本年政务，须按具体部众和地域复核。}] 明；朝廷与官员、边地政权与社群；边地接触区；不把族名或部名当固定疆界。朝贡名义、边境安全与港口贸易的利益使交涉持续发生。 它为后续册封、互市或海防安排留下文书与跨政权比较线索。 材料：《太祖实录》洪武七年条；《明史·外国传》；《朝鲜王朝实录》、琉球《历代宝案》或海外档案。限度：桥接条目只标示可回查的年度问题与材料链；实录有中央选择性，崇祯期尤其需要奏疏、地方、朝鲜及满文材料交叉。
  \item[\eventanchor{ML-1374-003}{1374（洪武七年）：沿海朝贡、海禁和港口管理的文书构成本年海上联系线索，禁令不能直接推作私人航行停止。}] 明；朝廷与官员；账本可辨空间；地点细节仍须回到所列材料。朝贡名义、边境安全与港口贸易的利益使交涉持续发生。 它为后续册封、互市或海防安排留下文书与跨政权比较线索。 材料：《太祖实录》洪武七年条；《明史·外国传》；《朝鲜王朝实录》、琉球《历代宝案》或海外档案。限度：桥接条目只标示可回查的年度问题与材料链；实录有中央选择性，崇祯期尤其需要奏疏、地方、朝鲜及满文材料交叉。
  \item[\eventanchor{ML-1375-001}{1375（洪武八年）：明开始发行新纸币“大明宝钞”}] 明；朝廷与官员；跨区域行政、资源或军政网络；不假定同步执行。继承、官僚协调或皇权运作的压力使问题进入朝廷程序。 它影响后续人事与政策议程；动机和派系标签不能由单一叙事裁定。 材料：《太祖实录》洪武八年条；《明史·本纪》及相关列传；诏令、奏疏或会典版本。限度：译写候选以编年材料定位；实录记载反映朝廷选择，崇祯期须另以奏疏、地方及敌对政权材料交叉。
  \item[\eventanchor{ML-1375-002}{1375（洪武八年）：由于费用和浪费，洪武皇帝停止了凤阳的建设；建设计划转移至南京}] 明；朝廷与官员；跨区域行政、资源或军政网络；不假定同步执行。继承、官僚协调或皇权运作的压力使问题进入朝廷程序。 它影响后续人事与政策议程；动机和派系标签不能由单一叙事裁定。 材料：《太祖实录》洪武八年条；《明史·本纪》及相关列传；诏令、奏疏或会典版本。限度：译写候选以编年材料定位；实录记载反映朝廷选择，崇祯期须另以奏疏、地方及敌对政权材料交叉。
  \item[\eventanchor{ML-1375-003}{1375（洪武八年）：皇室、功臣与中枢官僚的任免或案件材料构成本年权力重组线索，案狱叙事须避免照录罪名。}] 明；朝廷与官员；跨区域行政、资源或军政网络；不假定同步执行。继承、官僚协调或皇权运作的压力使问题进入朝廷程序。 它影响后续人事与政策议程；动机和派系标签不能由单一叙事裁定。 材料：《太祖实录》洪武八年条；《明史·本纪》及相关列传；诏令、奏疏或会典版本。限度：桥接条目只标示可回查的年度问题与材料链；实录有中央选择性，崇祯期尤其需要奏疏、地方、朝鲜及满文材料交叉。
  \item[\eventanchor{EM-1376-PROVINCIAL-COMMISSIONS}{1376（洪武九年）六月：废行中书省，分置承宣布政使司、都指挥使司、提刑按察使司}] 明；朝廷与官员、地方社会与精英；地方或省域；未具城镇者不补造路线。太祖防范行省集军政财权于一体，并需将新征服地区纳入可分察的行政链条 地方行政、军事和监察被分别归属三司；协调成本和中央派遣机制随之增加 材料：《太祖实录》洪武九年六月条；《明史·职官志》；《明会典·布政司》。限度：制度发布可靠；三司分置不等于地方权力已均衡运作
  \item[\eventanchor{ML-1376-001}{1376（洪武九年）三月：明军击败巴彦帖木儿}] 明；朝廷与官员、军队与卫所；跨区域行政、资源或军政网络；不假定同步执行。前期边防、地方武装或跨区域资源调动的压力推动了该行动。 它改变了后续调兵、补给或地方秩序的条件；战果和伤亡须分开核对。 材料：《太祖实录》洪武九年条；《明史·兵志》及相关列传；边镇、地方志或对方材料。限度：译写候选以编年材料定位；实录记载反映朝廷选择，崇祯期须另以奏疏、地方及敌对政权材料交叉。
  \item[\eventanchor{ML-1376-002}{1376（洪武九年）七月：明军再次击败巴彦帖木儿}] 明；朝廷与官员、军队与卫所；跨区域行政、资源或军政网络；不假定同步执行。前期边防、地方武装或跨区域资源调动的压力推动了该行动。 它改变了后续调兵、补给或地方秩序的条件；战果和伤亡须分开核对。 材料：《太祖实录》洪武九年条；《明史·兵志》及相关列传；边镇、地方志或对方材料。限度：译写候选以编年材料定位；实录记载反映朝廷选择，崇祯期须另以奏疏、地方及敌对政权材料交叉。
  \item[\eventanchor{ML-1376-004}{1376（洪武九年）：叶伯驹因批评皇帝重刑被饿死狱中}] 明；朝廷与官员；跨区域行政、资源或军政网络；不假定同步执行。继承、官僚协调或皇权运作的压力使问题进入朝廷程序。 它影响后续人事与政策议程；动机和派系标签不能由单一叙事裁定。 材料：《太祖实录》洪武九年条；《明史·本纪》及相关列传；诏令、奏疏或会典版本。限度：译写候选以编年材料定位；实录记载反映朝廷选择，崇祯期须另以奏疏、地方及敌对政权材料交叉。
  \item[\eventanchor{ML-1376-005}{1376（洪武九年）：洪武帝处决所有与“预印文书案”有关的官员}] 明；朝廷与官员；跨区域行政、资源或军政网络；不假定同步执行。继承、官僚协调或皇权运作的压力使问题进入朝廷程序。 它影响后续人事与政策议程；动机和派系标签不能由单一叙事裁定。 材料：《太祖实录》洪武九年条；《明史·本纪》及相关列传；诏令、奏疏或会典版本。限度：译写候选以编年材料定位；实录记载反映朝廷选择，崇祯期须另以奏疏、地方及敌对政权材料交叉。
  \item[\eventanchor{ML-1377-001}{1377（洪武十年）五月：明军入侵青海}] 明、后金（清）与辽东诸势力；朝廷与官员、军队与卫所、边地政权与社群；边地接触区；不把族名或部名当固定疆界。前期边防、地方武装或跨区域资源调动的压力推动了该行动。 它改变了后续调兵、补给或地方秩序的条件；战果和伤亡须分开核对。 材料：《太祖实录》洪武十年条；《明史·兵志》及相关列传；边镇、地方志或对方材料。限度：译写候选以编年材料定位；实录记载反映朝廷选择，崇祯期须另以奏疏、地方及敌对政权材料交叉。
  \item[\eventanchor{ML-1377-002}{1377（洪武十年）：南京宫殿竣工，定为“京师”}] 明；朝廷与官员、财政与商业行动者；跨区域行政、资源或军政网络；不假定同步执行。继承、官僚协调或皇权运作的压力使问题进入朝廷程序。 它影响后续人事与政策议程；动机和派系标签不能由单一叙事裁定。 材料：《太祖实录》洪武十年条；《明史·本纪》及相关列传；诏令、奏疏或会典版本。限度：译写候选以编年材料定位；实录记载反映朝廷选择，崇祯期须另以奏疏、地方及敌对政权材料交叉。
  \item[\eventanchor{ML-1377-003}{1377（洪武十年）：朝廷围绕卫所军户的编籍、屯田与调发整理本年军政文书，作为明初军事接收的可核查节点。}] 明；朝廷与官员、军队与卫所、财政与商业行动者；跨区域行政、资源或军政网络；不假定同步执行。前期边防、地方武装或跨区域资源调动的压力推动了该行动。 它改变了后续调兵、补给或地方秩序的条件；战果和伤亡须分开核对。 材料：《太祖实录》洪武十年条；《明史·兵志》及相关列传；边镇、地方志或对方材料。限度：桥接条目只标示可回查的年度问题与材料链；实录有中央选择性，崇祯期尤其需要奏疏、地方、朝鲜及满文材料交叉。
  \item[\eventanchor{ML-1378-001}{1378（洪武十一年）：吴冕叛乱：锦族叛乱}] 明与地方反抗、流动作战集团；朝廷与官员、军队与卫所、地方社会与精英、流动与反抗集团；地方或省域；未具城镇者不补造路线。前期边防、地方武装或跨区域资源调动的压力推动了该行动。 它改变了后续调兵、补给或地方秩序的条件；战果和伤亡须分开核对。 材料：《太祖实录》洪武十一年条；《明史·兵志》及相关列传；边镇、地方志或对方材料。限度：译写候选以编年材料定位；实录记载反映朝廷选择，崇祯期须另以奏疏、地方及敌对政权材料交叉。
  \item[\eventanchor{ML-1378-002}{1378（洪武十一年）：学校、科举、礼制或律令的颁行与讨论在本年材料中留下制度线索，规定与地方实践应分开处理。}] 明；朝廷与官员、地方社会与精英、宗教与文化行动者；地方或省域；未具城镇者不补造路线。制度、教育或知识传播的需要推动了相关行动或文本形成。 它提供制度和传播史的锚点，不能据此推断社会整体接受。 材料：《太祖实录》洪武十一年条；《明史·选举志》或《艺文志》；文集、碑刻或书目材料。限度：桥接条目只标示可回查的年度问题与材料链；实录有中央选择性，崇祯期尤其需要奏疏、地方、朝鲜及满文材料交叉。
  \item[\eventanchor{ML-1378-003}{1378（洪武十一年）：户部、布政司与里甲体系围绕户籍、田土和赋役的申报形成年度行政链，本条标记其可回查节点。}] 明；朝廷与官员、军队与卫所、财政与商业行动者；跨区域行政、资源或军政网络；不假定同步执行。财政征解、交通工程或货币支付的压力使相关措施进入政务。 其法定安排不等于地方执行，后续须以地域材料追踪负担与成效。 材料：《太祖实录》洪武十一年条；《明会典》相关条目；地方志、黄册/契约/河工材料。限度：桥接条目只标示可回查的年度问题与材料链；实录有中央选择性，崇祯期尤其需要奏疏、地方、朝鲜及满文材料交叉。
  \item[\eventanchor{ML-1379-001}{1379（洪武十二年）二月：明军在甘肃击败藏人}] 明；朝廷与官员、军队与卫所；跨区域行政、资源或军政网络；不假定同步执行。前期边防、地方武装或跨区域资源调动的压力推动了该行动。 它改变了后续调兵、补给或地方秩序的条件；战果和伤亡须分开核对。 材料：《太祖实录》洪武十二年条；《明史·兵志》及相关列传；边镇、地方志或对方材料。限度：译写候选以编年材料定位；实录记载反映朝廷选择，崇祯期须另以奏疏、地方及敌对政权材料交叉。
  \item[\eventanchor{ML-1379-002}{1379（洪武十二年）：占城向南京致敬}] 明与海上、朝贡相关政权；朝廷与官员；账本可辨空间；地点细节仍须回到所列材料。朝贡名义、边境安全与港口贸易的利益使交涉持续发生。 它为后续册封、互市或海防安排留下文书与跨政权比较线索。 材料：《太祖实录》洪武十二年条；《明史·外国传》；《朝鲜王朝实录》、琉球《历代宝案》或海外档案。限度：译写候选以编年材料定位；实录记载反映朝廷选择，崇祯期须另以奏疏、地方及敌对政权材料交叉。
  \item[\eventanchor{ML-1379-003}{1379（洪武十二年）：北元诸部、东北和西南的边报、招抚与防御命令持续进入本年政务，须按具体部众和地域复核。}] 明；朝廷与官员、边地政权与社群；边地接触区；不把族名或部名当固定疆界。朝贡名义、边境安全与港口贸易的利益使交涉持续发生。 它为后续册封、互市或海防安排留下文书与跨政权比较线索。 材料：《太祖实录》洪武十二年条；《明史·外国传》；《朝鲜王朝实录》、琉球《历代宝案》或海外档案。限度：桥接条目只标示可回查的年度问题与材料链；实录有中央选择性，崇祯期尤其需要奏疏、地方、朝鲜及满文材料交叉。
  \item[\eventanchor{EM-1380-HU-WEIYONG}{1380（洪武十三年）正月：胡惟庸案后中书省被废，皇帝直接统辖六部}] 明；朝廷与官员、地方社会与精英；地方或省域；未具城镇者不补造路线。太祖对宰相权力和官僚结党高度警惕；案件成为重组中枢权力的触发点 六部直接对皇帝负责；政务依赖皇帝裁断与后续内阁文书协调的矛盾加深 材料：《太祖实录》洪武十三年正月条；《明史·胡惟庸传》与《职官志》；《大诰》胡惟庸案条。限度：废相与制度变动可靠；胡惟庸“通倭通元”等指控须视为案狱叙事而非既定事实
  \item[\eventanchor{ML-1380-001}{1380（洪武十三年）：胡惟庸密谋刺杀洪武帝，被捕；随后的调查导致大约 15,000 人被处决}] 明；朝廷与官员；跨区域行政、资源或军政网络；不假定同步执行。继承、官僚协调或皇权运作的压力使问题进入朝廷程序。 它影响后续人事与政策议程；动机和派系标签不能由单一叙事裁定。 材料：《太祖实录》洪武十三年条；《明史·本纪》及相关列传；诏令、奏疏或会典版本。限度：译写候选以编年材料定位；实录记载反映朝廷选择，崇祯期须另以奏疏、地方及敌对政权材料交叉。
  \item[\eventanchor{ML-1380-002}{1380（洪武十三年）：“黄蜂巢”火箭炮是为明军制造的。}] 明；朝廷与官员、军队与卫所；跨区域行政、资源或军政网络；不假定同步执行。继承、官僚协调或皇权运作的压力使问题进入朝廷程序。 它影响后续人事与政策议程；动机和派系标签不能由单一叙事裁定。 材料：《太祖实录》洪武十三年条；《明史·本纪》及相关列传；诏令、奏疏或会典版本。限度：译写候选以编年材料定位；实录记载反映朝廷选择，崇祯期须另以奏疏、地方及敌对政权材料交叉。
  \item[\eventanchor{EM-1381-LIJIA-HUANGCE}{1381（洪武十四年）：朝廷推进黄册与里甲编制，登记户口、田土和役属}] 明与里甲户籍社会；朝廷与官员、财政与商业行动者；跨区域行政、资源或军政网络；不假定同步执行。新朝须将赋役、军户和土地责任落实到可追查的户等；元末流亡和战争破坏使登记困难 里甲成为赋役征派的基层框架；漏报、逃役和地方中介将持续塑造实际负担 材料：《太祖实录》洪武十四年户籍条；《明会典·户口》；徽州等地明代黄册残卷。限度：制度命令可靠；存世黄册仅反映部分地区与登记口径
  \item[\eventanchor{EM-1381-YUNNAN-CAMPAIGN}{1381（洪武十四年）：傅友德、蓝玉、沐英率军自川黔进攻云南}] 明、梁王与云南诸势力；朝廷与官员、军队与卫所；跨区域行政、资源或军政网络；不假定同步执行。梁王仍据云南并与北元保持联系；明廷欲控制西南交通、矿产和边疆名义 大规模军事接收开始；军屯、卫所和土司安排将深刻改变云南地方秩序 材料：《太祖实录》洪武十四年条；《明史·傅友德传》《沐英传》；《云南图经志书》云南府条。限度：出兵和主帅可靠；军额、伤亡与各族群立场不可照录明方奏报
  \item[\eventanchor{ML-1381-001}{1381（洪武十四年）十二月：明征服云南：明军占领曲靖}] 明；朝廷与官员、军队与卫所；跨区域行政、资源或军政网络；不假定同步执行。前期边防、地方武装或跨区域资源调动的压力推动了该行动。 它改变了后续调兵、补给或地方秩序的条件；战果和伤亡须分开核对。 材料：《太祖实录》洪武十四年条；《明史·兵志》及相关列传；边镇、地方志或对方材料。限度：译写候选以编年材料定位；实录记载反映朝廷选择，崇祯期须另以奏疏、地方及敌对政权材料交叉。
  \item[\eventanchor{ML-1382-001}{1382（洪武十五年）四月：明征服云南：明军征服云南}] 明；朝廷与官员、军队与卫所；跨区域行政、资源或军政网络；不假定同步执行。前期边防、地方武装或跨区域资源调动的压力推动了该行动。 它改变了后续调兵、补给或地方秩序的条件；战果和伤亡须分开核对。 材料：《太祖实录》洪武十五年条；《明史·兵志》及相关列传；边镇、地方志或对方材料。限度：译写候选以编年材料定位；实录记载反映朝廷选择，崇祯期须另以奏疏、地方及敌对政权材料交叉。
  \item[\eventanchor{EM-1383-JINYIWEI}{1383（洪武十六年）：设置锦衣卫，皇帝以亲军兼掌侦缉和诏狱}] 明；朝廷与官员、军队与卫所；跨区域行政、资源或军政网络；不假定同步执行。废相后皇帝强化对官僚和军政信息的直接掌握；大案审讯需要绕开常规司法链 亲军侦缉成为中枢权力工具之一；司法程序与政治案件的界线更趋紧张 材料：《太祖实录》洪武十六年条；《明史·职官志》；《明会典·锦衣卫》。限度：设卫可靠；后世将其职能倒投到洪武初年须谨慎
  \item[\eventanchor{ML-1383-001}{1383（洪武十六年）：沿海朝贡、海禁和港口管理的文书构成本年海上联系线索，禁令不能直接推作私人航行停止。}] 明；朝廷与官员；账本可辨空间；地点细节仍须回到所列材料。朝贡名义、边境安全与港口贸易的利益使交涉持续发生。 它为后续册封、互市或海防安排留下文书与跨政权比较线索。 材料：《太祖实录》洪武十六年条；《明史·外国传》；《朝鲜王朝实录》、琉球《历代宝案》或海外档案。限度：桥接条目只标示可回查的年度问题与材料链；实录有中央选择性，崇祯期尤其需要奏疏、地方、朝鲜及满文材料交叉。
  \item[\eventanchor{ML-1384-001}{1384（洪武十七年）四月：洪武皇帝将政府机构从皇宫迁至南京城墙外}] 明；朝廷与官员、地方社会与精英；地方或省域；未具城镇者不补造路线。继承、官僚协调或皇权运作的压力使问题进入朝廷程序。 它影响后续人事与政策议程；动机和派系标签不能由单一叙事裁定。 材料：《太祖实录》洪武十七年条；《明史·本纪》及相关列传；诏令、奏疏或会典版本。限度：译写候选以编年材料定位；实录记载反映朝廷选择，崇祯期须另以奏疏、地方及敌对政权材料交叉。
  \item[\eventanchor{ML-1384-002}{1384（洪武十七年）：皇室、功臣与中枢官僚的任免或案件材料构成本年权力重组线索，案狱叙事须避免照录罪名。}] 明；朝廷与官员；跨区域行政、资源或军政网络；不假定同步执行。继承、官僚协调或皇权运作的压力使问题进入朝廷程序。 它影响后续人事与政策议程；动机和派系标签不能由单一叙事裁定。 材料：《太祖实录》洪武十七年条；《明史·本纪》及相关列传；诏令、奏疏或会典版本。限度：桥接条目只标示可回查的年度问题与材料链；实录有中央选择性，崇祯期尤其需要奏疏、地方、朝鲜及满文材料交叉。
  \item[\eventanchor{EM-1385-DAGAO}{1385（洪武十八年）：《大诰》颁行，以案例和重刑语言约束官民}] 明；朝廷与官员；跨区域行政、资源或军政网络；不假定同步执行。太祖认为常律不足以威慑吏治和社会流动中的违法行为，并借案例塑造服从语言 法外重典进入基层教化与审判语境；法条、诏令和地方执行之间的落差须另行追踪 材料：《大诰》初编；《太祖实录》洪武十八年条；《明史·刑法志》。限度：文本与颁行可靠；规范性惩戒话语不能直接换算为实际犯罪率
  \item[\eventanchor{EM-1385-GUO-HUAN}{1385（洪武十八年）：郭桓案牵连户部及各地征粮官吏，朝廷严惩侵没赋税者}] 明与户部、地方官吏；朝廷与官员、军队与卫所、地方社会与精英、财政与商业行动者；地方或省域；未具城镇者不补造路线。江南赋粮征解经多层胥吏和官员转运，中央难以核对定额、耗损与实收 恐惧性稽核整顿财政链条但也冲击地方行政；账册真实度与官吏自保问题加剧 材料：《太祖实录》洪武十八年条；《明史·刑法志》；《大诰》郭桓案条。限度：案件与惩处可靠；被杀人数及“全国贪墨”规模高度依赖官方案狱口径
  \item[\eventanchor{ML-1385-001}{1385（洪武十八年）：吴冕叛乱：锦叛乱被击败}] 明与地方反抗、流动作战集团；朝廷与官员、军队与卫所、地方社会与精英、流动与反抗集团；地方或省域；未具城镇者不补造路线。前期边防、地方武装或跨区域资源调动的压力推动了该行动。 它改变了后续调兵、补给或地方秩序的条件；战果和伤亡须分开核对。 材料：《太祖实录》洪武十八年条；《明史·兵志》及相关列传；边镇、地方志或对方材料。限度：译写候选以编年材料定位；实录记载反映朝廷选择，崇祯期须另以奏疏、地方及敌对政权材料交叉。
  \item[\eventanchor{ML-1385-002}{1385（洪武十八年）：科举考试恢复}] 明；朝廷与官员、宗教与文化行动者；跨区域行政、资源或军政网络；不假定同步执行。制度、教育或知识传播的需要推动了相关行动或文本形成。 它提供制度和传播史的锚点，不能据此推断社会整体接受。 材料：《太祖实录》洪武十八年条；《明史·选举志》或《艺文志》；文集、碑刻或书目材料。限度：译写候选以编年材料定位；实录记载反映朝廷选择，崇祯期须另以奏疏、地方及敌对政权材料交叉。
  \item[\eventanchor{ML-1385-003}{1385（洪武十八年）：郭焕因贪污粮食700万担被执行死刑}] 明；朝廷与官员、财政与商业行动者；跨区域行政、资源或军政网络；不假定同步执行。继承、官僚协调或皇权运作的压力使问题进入朝廷程序。 它影响后续人事与政策议程；动机和派系标签不能由单一叙事裁定。 材料：《太祖实录》洪武十八年条；《明史·本纪》及相关列传；诏令、奏疏或会典版本。限度：译写候选以编年材料定位；实录记载反映朝廷选择，崇祯期须另以奏疏、地方及敌对政权材料交叉。
  \item[\eventanchor{ML-1386-001}{1386（洪武十九年）一月：明蒙毛战争：蒙毛起义军的司论法}] 明与地方反抗、流动作战集团；朝廷与官员、军队与卫所、地方社会与精英、流动与反抗集团；地方或省域；未具城镇者不补造路线。前期边防、地方武装或跨区域资源调动的压力推动了该行动。 它改变了后续调兵、补给或地方秩序的条件；战果和伤亡须分开核对。 材料：《太祖实录》洪武十九年条；《明史·兵志》及相关列传；边镇、地方志或对方材料。限度：译写候选以编年材料定位；实录记载反映朝廷选择，崇祯期须另以奏疏、地方及敌对政权材料交叉。
  \item[\eventanchor{ML-1386-002}{1386（洪武十九年）：朝廷围绕卫所军户的编籍、屯田与调发整理本年军政文书，作为明初军事接收的可核查节点。}] 明；朝廷与官员、军队与卫所、财政与商业行动者；跨区域行政、资源或军政网络；不假定同步执行。前期边防、地方武装或跨区域资源调动的压力推动了该行动。 它改变了后续调兵、补给或地方秩序的条件；战果和伤亡须分开核对。 材料：《太祖实录》洪武十九年条；《明史·兵志》及相关列传；边镇、地方志或对方材料。限度：桥接条目只标示可回查的年度问题与材料链；实录有中央选择性，崇祯期尤其需要奏疏、地方、朝鲜及满文材料交叉。
  \item[\eventanchor{EM-1387-FISHSCALE-REGISTERS}{1387（洪武二十年）：朝廷催成鱼鳞图册等田土登记以配合赋役核算}] 明与江南田土社会；朝廷与官员、地方社会与精英、财政与商业行动者；地方或省域；未具城镇者不补造路线。黄册登记户等后，国家需以地块和赋额关联田土责任；土地转移与隐匿造成核算压力 田土登记成为赋役分配的凭据；地方书手、业主和契约实践影响其可执行性 材料：《太祖实录》洪武二十年田土条；《明会典·田土》；徽州鱼鳞图册与相关契约。限度：命令与部分册籍可证；图册年代、抄补层次和亩数不等于实际占有
  \item[\eventanchor{EM-1387-NAGHACHU}{1387（洪武二十年）：冯胜等在辽东、松花江方向迫使纳哈出率部降明}] 明与纳哈出部；朝廷与官员、边地政权与社群；边地接触区；不把族名或部名当固定疆界。北元残部在辽东与松花江流域保持机动空间；明廷需稳住山海关以东通道 明在辽东增强卫所布置与招抚能力；东北诸部的朝贡和军事关系仍不断变动 材料：《太祖实录》洪武二十年条；《明史·冯胜传》与《北虏传》；《高丽史》辽东方面条。限度：降附事实可靠；部众数量和“归化”程度不能由明方献俘叙事推定
  \item[\eventanchor{ML-1387-001}{1387（洪武二十年）十月：明朝对乌里安凯的战役：那加楚向明朝军队投降}] 明与蒙古诸部；朝廷与官员、军队与卫所、边地政权与社群；边地接触区；不把族名或部名当固定疆界。前期边防、地方武装或跨区域资源调动的压力推动了该行动。 它改变了后续调兵、补给或地方秩序的条件；战果和伤亡须分开核对。 材料：《太祖实录》洪武二十年条；《明史·兵志》及相关列传；边镇、地方志或对方材料。限度：译写候选以编年材料定位；实录记载反映朝廷选择，崇祯期须另以奏疏、地方及敌对政权材料交叉。
  \item[\eventanchor{EM-1388-BUIR}{1388（洪武二十一年）：蓝玉军在捕鱼儿海击溃北元主力并俘获大批人员物资}] 明与北元；朝廷与官员、军队与卫所；跨区域行政、资源或军政网络；不假定同步执行。明军利用北元内部不稳和长距离奔袭寻求决定性战果 北元汗庭进一步分散；蓝玉军功骤升，也加深洪武朝对勋贵兵权的疑惧 材料：《太祖实录》洪武二十一年条；《明史·蓝玉传》与《北虏传》；17世纪《蒙古源流》北元汗统叙事。限度：胜利可靠；战果数字、元主脱逃经过和俘获规模应保留两方材料差异
  \item[\eventanchor{ML-1388-001}{1388（洪武二十一年）五月：贝尔湖之战：明军击败乌斯哈尔汗托古斯帖木儿}] 明；朝廷与官员、军队与卫所；跨区域行政、资源或军政网络；不假定同步执行。前期边防、地方武装或跨区域资源调动的压力推动了该行动。 它改变了后续调兵、补给或地方秩序的条件；战果和伤亡须分开核对。 材料：《太祖实录》洪武二十一年条；《明史·兵志》及相关列传；边镇、地方志或对方材料。限度：译写候选以编年材料定位；实录记载反映朝廷选择，崇祯期须另以奏疏、地方及敌对政权材料交叉。
  \item[\eventanchor{ML-1389-001}{1389（洪武二十二年）一月：明军在岳州击败彝族叛军}] 明与地方反抗、流动作战集团；朝廷与官员、军队与卫所、地方社会与精英、流动与反抗集团；地方或省域；未具城镇者不补造路线。前期边防、地方武装或跨区域资源调动的压力推动了该行动。 它改变了后续调兵、补给或地方秩序的条件；战果和伤亡须分开核对。 材料：《太祖实录》洪武二十二年条；《明史·兵志》及相关列传；边镇、地方志或对方材料。限度：译写候选以编年材料定位；实录记载反映朝廷选择，崇祯期须另以奏疏、地方及敌对政权材料交叉。
  \item[\eventanchor{ML-1389-002}{1389（洪武二十二年）十二月：明蒙战争：司伦法投降明朝}] 明；朝廷与官员、军队与卫所；跨区域行政、资源或军政网络；不假定同步执行。前期边防、地方武装或跨区域资源调动的压力推动了该行动。 它改变了后续调兵、补给或地方秩序的条件；战果和伤亡须分开核对。 材料：《太祖实录》洪武二十二年条；《明史·兵志》及相关列传；边镇、地方志或对方材料。限度：译写候选以编年材料定位；实录记载反映朝廷选择，崇祯期须另以奏疏、地方及敌对政权材料交叉。
  \item[\eventanchor{EM-1390-LAN-YU}{1390（洪武二十三年）：蓝玉以谋反罪被诛，功臣集团再受清洗}] 明；朝廷与官员；跨区域行政、资源或军政网络；不假定同步执行。北伐功臣掌握军中网络且太祖持续防范结党；案狱成为重排军政人事的渠道 开国武将集团受到重创；皇太孙继承安排更依赖文官和制度性控制 材料：《太祖实录》洪武二十三年条；《明史·蓝玉传》；《大诰》相关案件。限度：处决与案名可靠；谋反细节主要来自后修实录和官方传记，难以独立裁定
  \item[\eventanchor{ML-1390-001}{1390（洪武二十三年）四月：纳伊尔布哈和耀珠向明军投降}] 明；朝廷与官员、军队与卫所；跨区域行政、资源或军政网络；不假定同步执行。前期边防、地方武装或跨区域资源调动的压力推动了该行动。 它改变了后续调兵、补给或地方秩序的条件；战果和伤亡须分开核对。 材料：《太祖实录》洪武二十三年条；《明史·兵志》及相关列传；边镇、地方志或对方材料。限度：译写候选以编年材料定位；实录记载反映朝廷选择，崇祯期须另以奏疏、地方及敌对政权材料交叉。
  \item[\eventanchor{ML-1391-001}{1391（洪武二十四年）五月：阿贾施里叛乱并被镇压}] 明与地方反抗、流动作战集团；朝廷与官员、军队与卫所、地方社会与精英、流动与反抗集团；地方或省域；未具城镇者不补造路线。继承、官僚协调或皇权运作的压力使问题进入朝廷程序。 它影响后续人事与政策议程；动机和派系标签不能由单一叙事裁定。 材料：《太祖实录》洪武二十四年条；《明史·本纪》及相关列传；诏令、奏疏或会典版本。限度：译写候选以编年材料定位；实录记载反映朝廷选择，崇祯期须另以奏疏、地方及敌对政权材料交叉。
  \item[\eventanchor{ML-1391-002}{1391（洪武二十四年）：明军短暂占领哈密并撤退}] 明；朝廷与官员、军队与卫所、边地政权与社群；边地接触区；不把族名或部名当固定疆界。前期边防、地方武装或跨区域资源调动的压力推动了该行动。 它改变了后续调兵、补给或地方秩序的条件；战果和伤亡须分开核对。 材料：《太祖实录》洪武二十四年条；《明史·兵志》及相关列传；边镇、地方志或对方材料。限度：译写候选以编年材料定位；实录记载反映朝廷选择，崇祯期须另以奏疏、地方及敌对政权材料交叉。
  \item[\eventanchor{EM-1392-HEIR-CHANGE}{1392（洪武二十五年）四月：皇太子朱标去世，太祖立朱允炆为皇太孙}] 明；朝廷与官员；跨区域行政、资源或军政网络；不假定同步执行。朱标早逝打破原有继承顺序；太祖需在皇孙与成年藩王之间确立可被承认的名分 建文君臣日后削藩拥有继承政治背景；燕王等边王的资源与疑虑同时上升 材料：《太祖实录》洪武二十五年四月条；《明史·懿文太子传》与《惠帝本纪》；南京懿文太子陵考古简报与碑刻著录。限度：继承变更可靠；对朱标能力和诸王态度的评价多为后出叙事
  \item[\eventanchor{ML-1392-001}{1392（洪武二十五年）八月5日：李松桂废黜王耀，成为朝鲜太祖；高句丽就此结束}] 明与海上、朝贡相关政权；朝廷与官员；账本可辨空间；地点细节仍须回到所列材料。朝贡名义、边境安全与港口贸易的利益使交涉持续发生。 它为后续册封、互市或海防安排留下文书与跨政权比较线索。 材料：《太祖实录》洪武二十五年条；《明史·外国传》；《朝鲜王朝实录》、琉球《历代宝案》或海外档案。限度：译写候选以编年材料定位；实录记载反映朝廷选择，崇祯期须另以奏疏、地方及敌对政权材料交叉。
  \item[\eventanchor{EM-1393-BLUE-JADE-PURGE}{1393（洪武二十六年）：蓝玉案牵连扩大，勋贵与官员大批获罪}] 明；朝廷与官员；跨区域行政、资源或军政网络；不假定同步执行。太祖以案狱消除潜在军事和官僚自主网络，并为皇太孙预作权力清场 朝廷人事更加依附皇权；以高压保障继承也削弱了可制衡藩王的开国武臣 材料：《太祖实录》洪武二十六年条；《明史·功臣世表》与相关列传；《大诰三编》。限度：清洗持续及其政治影响可靠；受牵连人数和罪证须避免采纳单一官方总数
  \item[\eventanchor{ML-1393-001}{1393（洪武二十六年）：明军攻陷哈密}] 明；朝廷与官员、军队与卫所、边地政权与社群；边地接触区；不把族名或部名当固定疆界。前期边防、地方武装或跨区域资源调动的压力推动了该行动。 它改变了后续调兵、补给或地方秩序的条件；战果和伤亡须分开核对。 材料：《太祖实录》洪武二十六年条；《明史·兵志》及相关列传；边镇、地方志或对方材料。限度：译写候选以编年材料定位；实录记载反映朝廷选择，崇祯期须另以奏疏、地方及敌对政权材料交叉。
  \item[\eventanchor{ML-1394-001}{1394（洪武二十七年）：明朝鲜朝贡关系正常化}] 明与海上、朝贡相关政权；朝廷与官员；账本可辨空间；地点细节仍须回到所列材料。朝贡名义、边境安全与港口贸易的利益使交涉持续发生。 它为后续册封、互市或海防安排留下文书与跨政权比较线索。 材料：《太祖实录》洪武二十七年条；《明史·外国传》；《朝鲜王朝实录》、琉球《历代宝案》或海外档案。限度：译写候选以编年材料定位；实录记载反映朝廷选择，崇祯期须另以奏疏、地方及敌对政权材料交叉。
  \item[\eventanchor{ML-1394-002}{1394（洪武二十七年）：户部、布政司与里甲体系围绕户籍、田土和赋役的申报形成年度行政链，本条标记其可回查节点。}] 明；朝廷与官员、军队与卫所、财政与商业行动者；跨区域行政、资源或军政网络；不假定同步执行。财政征解、交通工程或货币支付的压力使相关措施进入政务。 其法定安排不等于地方执行，后续须以地域材料追踪负担与成效。 材料：《太祖实录》洪武二十七年条；《明会典》相关条目；地方志、黄册/契约/河工材料。限度：桥接条目只标示可回查的年度问题与材料链；实录有中央选择性，崇祯期尤其需要奏疏、地方、朝鲜及满文材料交叉。
  \item[\eventanchor{ML-1395-001}{1395（洪武二十八年）：户部、布政司与里甲体系围绕户籍、田土和赋役的申报形成年度行政链，本条标记其可回查节点。}] 明；朝廷与官员、军队与卫所、财政与商业行动者；跨区域行政、资源或军政网络；不假定同步执行。财政征解、交通工程或货币支付的压力使相关措施进入政务。 其法定安排不等于地方执行，后续须以地域材料追踪负担与成效。 材料：《太祖实录》洪武二十八年条；《明会典》相关条目；地方志、黄册/契约/河工材料。限度：桥接条目只标示可回查的年度问题与材料链；实录有中央选择性，崇祯期尤其需要奏疏、地方、朝鲜及满文材料交叉。
  \item[\eventanchor{ML-1395-002}{1395（洪武二十八年）：北元诸部、东北和西南的边报、招抚与防御命令持续进入本年政务，须按具体部众和地域复核。}] 明；朝廷与官员、边地政权与社群；边地接触区；不把族名或部名当固定疆界。朝贡名义、边境安全与港口贸易的利益使交涉持续发生。 它为后续册封、互市或海防安排留下文书与跨政权比较线索。 材料：《太祖实录》洪武二十八年条；《明史·外国传》；《朝鲜王朝实录》、琉球《历代宝案》或海外档案。限度：桥接条目只标示可回查的年度问题与材料链；实录有中央选择性，崇祯期尤其需要奏疏、地方、朝鲜及满文材料交叉。
  \item[\eventanchor{ML-1396-001}{1396（洪武二十九年）四月：明军击败勃林帖木儿}] 明；朝廷与官员、军队与卫所；跨区域行政、资源或军政网络；不假定同步执行。前期边防、地方武装或跨区域资源调动的压力推动了该行动。 它改变了后续调兵、补给或地方秩序的条件；战果和伤亡须分开核对。 材料：《太祖实录》洪武二十九年条；《明史·兵志》及相关列传；边镇、地方志或对方材料。限度：译写候选以编年材料定位；实录记载反映朝廷选择，崇祯期须另以奏疏、地方及敌对政权材料交叉。
  \item[\eventanchor{ML-1396-002}{1396（洪武二十九年）：学校、科举、礼制或律令的颁行与讨论在本年材料中留下制度线索，规定与地方实践应分开处理。}] 明；朝廷与官员、地方社会与精英、宗教与文化行动者；地方或省域；未具城镇者不补造路线。制度、教育或知识传播的需要推动了相关行动或文本形成。 它提供制度和传播史的锚点，不能据此推断社会整体接受。 材料：《太祖实录》洪武二十九年条；《明史·选举志》或《艺文志》；文集、碑刻或书目材料。限度：桥接条目只标示可回查的年度问题与材料链；实录有中央选择性，崇祯期尤其需要奏疏、地方、朝鲜及满文材料交叉。
  \item[\eventanchor{EM-1397-DAMING-LAW}{1397（洪武三十年）：定颁《大明律》终本，形成以六律为纲的成文法框架}] 明；朝廷与官员；跨区域行政、资源或军政网络；不假定同步执行。洪武前期反复修律与《大诰》重典并行，中央需要稳定、可传抄的法典文本 终本成为明代司法基本参照；诏令、条例和诏狱仍会改变实际司法面貌 材料：《大明律》洪武三十年刊本系统；《太祖实录》洪武三十年条；《明史·刑法志》。限度：终本年代可靠；法典条文不能证明州县审判或社会遵从一致
  \item[\eventanchor{ML-1397-001}{1397（洪武三十年）十月：林宽叛乱：锦族叛乱被击败}] 明与地方反抗、流动作战集团；朝廷与官员、军队与卫所、地方社会与精英、流动与反抗集团；地方或省域；未具城镇者不补造路线。前期边防、地方武装或跨区域资源调动的压力推动了该行动。 它改变了后续调兵、补给或地方秩序的条件；战果和伤亡须分开核对。 材料：《太祖实录》洪武三十年条；《明史·兵志》及相关列传；边镇、地方志或对方材料。限度：译写候选以编年材料定位；实录记载反映朝廷选择，崇祯期须另以奏疏、地方及敌对政权材料交叉。
  \item[\eventanchor{ML-1397-002}{1397（洪武三十年）十二月：明蒙毛干预：斯伦法被废黜并请求明援助以恢复其权力}] 明；朝廷与官员；跨区域行政、资源或军政网络；不假定同步执行。继承、官僚协调或皇权运作的压力使问题进入朝廷程序。 它影响后续人事与政策议程；动机和派系标签不能由单一叙事裁定。 材料：《太祖实录》洪武三十年条；《明史·本纪》及相关列传；诏令、奏疏或会典版本。限度：译写候选以编年材料定位；实录记载反映朝廷选择，崇祯期须另以奏疏、地方及敌对政权材料交叉。
  \item[\eventanchor{EM-1398-HONGWU-DEATH}{1398（洪武三十一年）闰五月：太祖去世，皇太孙朱允炆即位为建文帝}] 明；朝廷与官员；跨区域行政、资源或军政网络；不假定同步执行。太祖已以皇太孙继承并清洗部分功臣，但成年藩王仍掌边镇卫所资源 建文新政面对削藩与守成的选择；中央—燕王冲突迅速制度化 材料：《太祖实录》洪武三十一年闰五月条；《明史·惠帝本纪》；孝陵碑刻与陵寝考古资料。限度：继位日期可靠；遗诏与宫廷部署的细节受建文、永乐两朝修史影响
  \item[\eventanchor{ML-1398-001}{1398（洪武三十一年）一月：明蒙毛干预：斯伦法恢复权力}] 明；朝廷与官员；跨区域行政、资源或军政网络；不假定同步执行。继承、官僚协调或皇权运作的压力使问题进入朝廷程序。 它影响后续人事与政策议程；动机和派系标签不能由单一叙事裁定。 材料：《太祖实录》洪武三十一年条；《明史·本纪》及相关列传；诏令、奏疏或会典版本。限度：译写候选以编年材料定位；实录记载反映朝廷选择，崇祯期须另以奏疏、地方及敌对政权材料交叉。
  \item[\eventanchor{ML-1398-002}{1398（洪武三十一年）五月24日：洪武帝生病了}] 明；朝廷与官员；跨区域行政、资源或军政网络；不假定同步执行。继承、官僚协调或皇权运作的压力使问题进入朝廷程序。 它影响后续人事与政策议程；动机和派系标签不能由单一叙事裁定。 材料：《太祖实录》洪武三十一年条；《明史·本纪》及相关列传；诏令、奏疏或会典版本。限度：译写候选以编年材料定位；实录记载反映朝廷选择，崇祯期须另以奏疏、地方及敌对政权材料交叉。
  \item[\eventanchor{ML-1398-003}{1398（洪武三十一年）六月24日：洪武帝驾崩}] 明；朝廷与官员；跨区域行政、资源或军政网络；不假定同步执行。继承、官僚协调或皇权运作的压力使问题进入朝廷程序。 它影响后续人事与政策议程；动机和派系标签不能由单一叙事裁定。 材料：《太祖实录》洪武三十一年条；《明史·本纪》及相关列传；诏令、奏疏或会典版本。限度：译写候选以编年材料定位；实录记载反映朝廷选择，崇祯期须另以奏疏、地方及敌对政权材料交叉。
  \item[\eventanchor{ML-1398-004}{1398（洪武三十一年）六月30日：朱允炆登基为建文帝}] 明；朝廷与官员；跨区域行政、资源或军政网络；不假定同步执行。继承、官僚协调或皇权运作的压力使问题进入朝廷程序。 它影响后续人事与政策议程；动机和派系标签不能由单一叙事裁定。 材料：《太祖实录》洪武三十一年条；《明史·本纪》及相关列传；诏令、奏疏或会典版本。限度：译写候选以编年材料定位；实录记载反映朝廷选择，崇祯期须另以奏疏、地方及敌对政权材料交叉。
  \item[\eventanchor{ML-1398-005}{1398（洪武三十一年）：建文帝灭朱珪、朱勃、朱辅、朱翩诸侯国}] 明；朝廷与官员；跨区域行政、资源或军政网络；不假定同步执行。继承、官僚协调或皇权运作的压力使问题进入朝廷程序。 它影响后续人事与政策议程；动机和派系标签不能由单一叙事裁定。 材料：《太祖实录》洪武三十一年条；《明史·本纪》及相关列传；诏令、奏疏或会典版本。限度：译写候选以编年材料定位；实录记载反映朝廷选择，崇祯期须另以奏疏、地方及敌对政权材料交叉。
  \item[\eventanchor{ML-1398-006}{1398（洪武三十一年）：中国最后一次活人祭祀记录}] 明；朝廷与官员；跨区域行政、资源或军政网络；不假定同步执行。继承、官僚协调或皇权运作的压力使问题进入朝廷程序。 它影响后续人事与政策议程；动机和派系标签不能由单一叙事裁定。 材料：《太祖实录》洪武三十一年条；《明史·本纪》及相关列传；诏令、奏疏或会典版本。限度：译写候选以编年材料定位；实录记载反映朝廷选择，崇祯期须另以奏疏、地方及敌对政权材料交叉。
\end{description}
